\chapter*{General Introduction}
\addcontentsline{toc}{chapter}{General Introduction}
\begin{spacing}{1.2}
%==============================================================================

Writing a good computer-science project report \cite{SFAXI2015} follows a few
universal rules. Everyone writes with their own voice and signature, but some
conventions remain shared \cite{Latex}.\\

\textbf{The Table of Contents} is the first thing a reviewer reads. It should
be:
\begin{itemize}
  \item Detailed enough\footnote{But not too detailed} -- three numbering levels
        usually suffice;
  \item Balanced across chapters, except for the Introduction and the
        Conclusion which are naturally shorter than the others;
  \item Built from descriptive titles. Avoid generic ``Design''; prefer
        ``Design of the management application for $\ldots$''. Slightly longer
        titles are better than an impersonal summary.\\
\end{itemize}

\textbf{An introduction} should be written as well-crafted paragraphs.
It usually has four parts:
\begin{enumerate}
  \item \textbf{Context} -- the general domain of your application (web, BI,
        enterprise software, etc.);
  \item \textbf{Problem statement} -- what needs in this context motivate your
        project;
  \item \textbf{Contribution} -- a short description of your application,
        without going into implementation details. The introduction introduces
        the work, it does not summarize it;
  \item \textbf{Report outline} -- the different chapters and what they cover.
        These items do not need to be numbered; prefer well-connected
        paragraphs.
\end{enumerate}

\end{spacing}
